\section{A genuine cuspidal-support construction in the ordinary split branch}
\label{sec:cubic-hsieh-support}

The following result removes one premise from a particular
Type I construction: nonunitness of its imprimitive
$p$-adic $L$-function can be arranged by the level.
The assertion concerns cuspidal Hecke support; its scope
at the end distinguishes this from an irreducible seed.
The cited Hida and nonvanishing theorems are retained
published inputs, not consequences of Theorem A.

\begin{proposition}[AC-K29-Hsieh-aux: auxiliary level forces nonzero cuspidal support]
\label{AC-K29-Hsieh-aux}
Let $K$ be an imaginary quadratic field, let $p>3$
split in $K$, and suppose $p\nmid h_K$.
Let $\beta$ be a finite-order character of $G_K$ of
order prime to $p$, unramified at both primes over $p$.
Use Hsieh's ordinary two-variable branch
\[
 \Lambda=W_p[\beta][[\Gamma_K]],\qquad
                 \Psi=\beta\varepsilon_\Lambda .          \tag{HS.1}
\]
Retain the ordinary Eisenstein construction and control
theorems of Hsieh, \emph{Eisenstein congruence on unitary
groups and Iwasawa main conjectures for CM fields},
\S\S4--7, with the simultaneous nonvanishing input
as completed by Burungale--He--Tian--Ye,
arXiv:2508.19706v2, Theorems 6.2 and 6.5.

One can choose its prime-to-$p$ auxiliary level so that
the local ordinary \emph{cuspidal} Hecke algebra
$R=h^{\rm ord}_{E}$ is nonzero and finite,
reduced and torsion-free over $\Lambda$, and
\[
 R/I^{S}(\Psi)\twoheadrightarrow
 \Lambda/\bigl(L_p^{S}(-1,\Psi^+\tau_{K/\mathbf Q})
                         L_p^{S}(\Psi,\Sigma)\bigr)
                  \ne0 .                                  \tag{HS.2}
\]
Every classical point of this localized component has
residual semisimplification
\[
                 \bar\beta^{-c}\bar\kappa_p
                       \oplus1\oplus
                         \bar\beta\bar\kappa_p^{-1}.       \tag{HS.3}
\]
Classical ordinary cuspidal points exist in this component.
No nonzero Selmer class or special-value divisibility
is a premise of this construction.
\end{proposition}
\begin{proof}
Let $K_0$ be the compositum of the field cut out by
$\beta$ with $K(\mu_p)$.  Its degree over $K$ is prime
to $p$, so it is disjoint from the pro-$p$ extension
$K_\infty/K$ used in (HS.1).  It satisfies the finite
branch-field hypothesis of the cited sources.
At every nonsplit prime in the conductor of $\beta$,
some inertia value is a nontrivial root of unity of
order prime to $p$.  Its difference from $1$ is a
$p$-adic unit.  Thus the invariant in Hsieh (7.11)
satisfies $\mu_p(\beta)=0$.
The conjugate-prime branch condition holds, and the
integer called $a$ there can be chosen to be $0$.
In particular $a\not\equiv2\pmod{p-1}$.
The repaired nonvanishing theorem applies since
$p\nmid6h_K[K_0:K]$.
Fix the initial Fourier--Jacobi triple and its auxiliary
nonvanishing prime supplied by these inputs before
choosing $q$, and include their finite bad places in
the excluded set below.

Choose a rational prime $q$ completely split in the
normal closure of $K_0/\mathbf Q$, outside the fixed
bad primes, including the conductor, $p$, the
discriminant, and the fixed Hermitian lattice's bad
primes.  Existence requires no estimate for the least
such prime.  For completeness it follows from the
elementary prime-divisor argument: the minimal
polynomial of an integral primitive element of this
normal closure has prime divisors outside every fixed
finite set.  Indeed if $f(a)\ne0$ and $M$ is divisible
by that set, then
$f(a+Mf(a)t)/f(a)$ is an integer congruent to $1$
modulo $M$ and has absolute value greater than $1$
for large $t$.  A prime divisor outside the set gives
a linear factor of $f$; excluding its discriminant
and index, normality makes the corresponding Frobenius
trivial, so the prime is completely split.

Require $q$ to divide the split-prime product $s$ in
Hsieh \S5.1.1(a1), printed p.~52.
This is compatible with its fixed Hermitian-lattice
condition (h3) by our choice of $q$.
Choose the remaining level sufficiently deep to
satisfy (a2)--(a3).
This prescribed extra split prime does not invalidate
primitivity: in the complete proof of Proposition 7.13,
printed pp.~87--88, the level is chosen after the
initial Fourier--Jacobi triple.  Every added split
prime only adds a finite set of excluded values of
the auxiliary anticyclotomic character $\nu$.
The simultaneous nonvanishing set is infinite, while
each such local exclusion is finite, by the same
nontrivial Frobenius argument in that proof.
The replacement simultaneous nonvanishing theorem
just cited supplies exactly this infinite set.
Thus its chosen Fourier--Jacobi functional remains
a unit at this prescribed enlarged level.

For either $w\mid q$, the imprimitive-factor formula
in Hsieh \S6.5, immediately before (6.14), printed
p.~75, includes
\[
               1-\Psi(\operatorname{Frob}_w)
                        \in\mathfrak m_\Lambda.             \tag{HS.4}
\]
Indeed $\beta(\operatorname{Frob}_w)=1$ and every
group-like element of $\Gamma_K$ reduces to $1$
modulo $\mathfrak m_\Lambda$.
The primitive Katz function is in $\Lambda$.
The imprimitive Deligne--Ribet function is also in
$\Lambda$, as stated at the same locator; $S$ is
nonempty.  Consequently their displayed product
belongs to $\mathfrak m_\Lambda$ and is a nonunit.

The full proof of Hsieh Theorem 7.18, printed
pp.~90--91, now applies.  It subtracts the constant
term times an integral ordinary form and obtains a
cuspidal form with a unit Fourier--Jacobi functional.
This gives (HS.2).  Definition 7.19 gives the
nonzero finite reduced torsion-free localized algebra.
Its finite map to $\operatorname{Spec}\Lambda$ is
surjective: its annihilator over $\Lambda$ is zero,
and finite modules have closed support.
Proposition 7.17 and the ordinary control theorem
therefore supply classical cuspidal points above
classical weights.  Finally the complete construction
in \S8.3, equation (8.2), printed pp.~95--96,
identifies the residual pseudocharacter as the sum
in (HS.3).  This proves the assertions without using
a Selmer group to force (HS.4).
\end{proof}

\begin{corollary}[AC-K29-Hsieh-branch: exact cubic specialization and range]
\label{AC-K29-Hsieh-branch}
Take the fixed cubic anticyclotomic $\chi$, $j=1,2$,
and the rational character $\gamma=\psi\kappa_p^e$
of the actual source, with $\psi$ of order prime to
$p$, $\psi^2|_{I_p}=1$, and
$p-1\nmid2e$.  Suppose $p>3$ splits in $K$, and
\[
       \bar\psi|_{I_p}=\bar\kappa_p^{\,1-e}|_{I_p}.
                                                               \tag{HS.5}
\]
Then AC-K29-Hsieh-aux applies to
\[
                  \beta=\chi^j\psi^{-1}\omega_p^{1-e},       \tag{HS.6}
\]
and (HS.3) is exactly
$\chi^j\bar\gamma\oplus1\oplus
\chi^j\bar\gamma^{-1}$.
For $e>1$, (HS.5) entails
\[
                         p-1\mid2(e-1),\qquad p\le2e-1.     \tag{HS.7}
\]
The equality $e\equiv1\pmod{p-1}$ follows only if
$\psi$ itself is unramified at $p$.
\end{corollary}
\begin{proof}
The cubic character is unramified at this split
coefficient prime.  The two local restrictions of
the rational $\psi$ and $\omega_p$ at the primes of
$K$ above $p$ agree.  Thus (HS.5) makes (HS.6)
unramified at both primes.  Its order is prime to
$p$, so the proposition applies.
Conjugation inverts $\chi$ and fixes the rational
factors; direct substitution gives (HS.3).
Squaring (HS.5), using $\bar\psi^2=1$ and the full
order $p-1$ of the cyclotomic inertia character,
proves (HS.7).  The last assertion follows by
removing $\bar\psi$ in (HS.5).
\end{proof}

\paragraph{Exact remaining interface.}
The actual target specialization
$\Psi=\chi^j\gamma^{-1}\kappa_p$ has parallel Tate
infinity type $(1-e,1-e)$; for $e\ge1$ it is outside
the regular classical cone used in Hsieh \S8.4,
printed pp.~97--98.  The construction instead uses
classical points in its residual branch.
Hsieh \S8.3 attaches a \emph{semisimple}
three-dimensional representation to those cuspidal
unitary points.  The proposition does not prove
that one of these representations is absolutely
irreducible, satisfies the prescribed seed weights
and away-from-$p$ types, or propagates to all the
modified deformation points.
It gives no uniform bound for the auxiliary level.
These are the additional actual seed and level
requirements in AC-K29-TypeI-source.
For other branches or inert coefficient primes the
stated construction has not been supplied here.

\paragraph{Geometric dictionary.}
The branch is a rank-one $p$-adic local system; the
three terms in (HS.3) give its residual unitary
deformation datum.  The new conclusion is nonzero
cuspidal Hecke support at an explicitly permitted
auxiliary level.  The finite-cover terminology has
the previously audited SGA~1, Expos\'e~V, \S\S4 and~7
meaning only; the Hecke-support and primitivity
statements come from the cited automorphic inputs.
